\chapter{Functional Requirements}
\label{ch:functional-requirements}

This chapter defines the functional requirements that the Sentinel system
must satisfy. Requirements are described functionally, independent of the
specific hardware components used to realise them; component selection and
justification are deferred to dedicated design chapters later in this
report.

\section{Climbing Mechanics and Slack Management}
\label{sec:fr-slack-management}

The system must automatically take in rope as the climber ascends, so
that the amount of slack in the system remains within a safe, comfortable
range without manual intervention. To do so, the system must reliably
distinguish between two dynamically different situations that both
involve rope moving through the mechanism:

\begin{itemize}
    \item \textbf{Normal ascent / rope pull-in}, where the climber is
        moving upward and the rope must be taken in at a corresponding,
        comparatively gradual rate. The device also has to be able to allow
        rope taking during climbing without blocking the system.
    \item \textbf{A fall event}, characterised by a sudden increase in
        rope acceleration and a rapid increase in rope
        tension, which must be distinguished from normal climbing
        dynamics (e.g.\ clipping, resting on the rope, or bursts of quick
        movement).
\end{itemize}

This distinction must be made using onboard rope-movement and
tension-sensing information, processed in real time, and must be robust
enough to avoid both missed fall detections and false-positive brake
triggers during ordinary climbing.

\section{Post-Fall State Detection}
\label{sec:fr-post-fall}

After the system has arrested a fall by engaging the brake, it must be
able to detect when the climber has resumed a valid, stable climbing
position (e.g.\ has regained a hold and is no longer hanging dynamically
on the rope). Once this state is detected, the system must release the
brake and resume normal slack-management behaviour automatically, without
requiring a manual reset from the climber or a belay partner, unless a
manual reset is explicitly designed as a deliberate safety fallback. Also
this has to be able to do by pressing a button in the remote control.

\section{Operational Modes and Remote Control}
\label{sec:fr-modes}

The system must support at least two distinct operational modes:

\begin{itemize}
    \item \textbf{Climbing Mode}: the default mode described in
        \cref{sec:fr-slack-management} and \cref{sec:fr-post-fall}, in
        which the system manages slack during ascent and arrests falls.
    \item \textbf{Rappel Mode}: a mode that allows a controlled,
        continuous descent, with the braking behaviour adapted for a
        deliberate downward motion rather than treating it as a fall
        event.
\end{itemize}

Switching between these modes must be possible via a wireless remote
control, allowing the operator (climber or belay supervisor) to change
modes without physically accessing the device mounted on the anchor point.

The rappel will be done by maintaining a downwards button in the remote controll.
There must be an other button in the remote controll to pull the rope.

\section{Battery Management and Safety Protocols}
\label{sec:fr-battery}

\subsection{Battery Level Indication}

The system must visually and clearly indicate the current battery charge
level at all times during operation, so that the user can make an
informed decision about whether it is safe to begin or continue climbing.

\subsection{Low-Battery Restrictions}

The system must enforce the following safety protocols based on battery
state:

\begin{itemize}
    \item If the battery level is low at start-up, the system must
        prevent Climbing Mode from being initiated.
    \item If the battery level drops to a low or critical level
        \emph{during} an active climbing session, the system must
        restrict its functionality strictly to Rappel/Abseil Mode,
        allowing the climber currently on the wall to be safely lowered
        rather than left in an unsupported state. The system must do a clear 
        sonoric advice to announce the mode switch.
\end{itemize}

\section{Fail-Safe Mechanical Override}
\label{sec:fr-failsafe}

This requirement is considered \textbf{critical} to the safety case of
the system.

\subsection{Automatic Rope Lock on Power Loss}

In the event of a power loss or an electronic/software failure, the rope
must lock automatically by default. The mechanical design must guarantee
this fail-safe behaviour independent of the control software -- that is,
the default, unpowered, or faulted state of the mechanism must be
``locked'', not ``free-running''.

\subsection{Manual Override Lever}

A manual physical lever must allow the climber (or a belay supervisor) to
override the automatic lock and execute a controlled descent by hand,
even with no power available. This requires evaluating a passive,
GriGri-like mechanical behaviour that functions correctly without any
electronic assistance, ensuring that a power or electronics failure never
leaves a climber stranded with no means of controlled descent.

\section{Non-Functional Requirements}
\label{sec:non-functional-requirements}

Beyond the functional behaviour described above, the system must also
satisfy a number of non-functional requirements. These are not yet
exhaustive and are expected to be expanded as the design matures; initial
candidates include:

\begin{itemize}
    \item \textbf{Reliability and safety-critical performance}: given
        that a missed or delayed fall detection has direct safety
        consequences, the system's sensing and braking response must
        meet a bounded, worst-case reaction latency, and must degrade
        towards the fail-safe locked state (\cref{sec:fr-failsafe})
        rather than towards an unsafe state under any fault condition.
    \item \textbf{Real-time responsiveness}: sensor acquisition, state
        estimation, and brake actuation must run within firm real-time
        deadlines, since the usefulness of the system depends on
        reacting fast enough to arrest a fall before excessive slack
        develops.
    \item \textbf{Usability and ergonomics}: the device must be
        operable by a climber with no more training than is required for
        a conventional assisted-braking belay device, with clear,
        unambiguous feedback (e.g.\ battery level, operating mode) that
        does not require reading a manual mid-climb.
\end{itemize}
